% =====================================================================
%  Journal article skeleton for TAG-MTL-DR.
%  Every number comes from results_macros.tex and the generated tables
%  (analysis.py). Text in [brackets] is to be written once results exist;
%  do not state a claim the tables do not support.
%  Build: set \resultsdir, then pdflatex + bibtex.
% =====================================================================
\documentclass[11pt,a4paper]{article}
\usepackage[margin=2.2cm]{geometry}
\usepackage{booktabs,graphicx,amsmath,amssymb,xcolor,hyperref,cleveref}
\newcommand{\resultsdir}{../runs/main/paper}
\input{\resultsdir/results_macros.tex}
\newcommand{\todo}[1]{\textcolor{red}{[#1]}}

\title{Topology-Aware Graph and Multi-Task Learning for Safety-Oriented
Diabetic Retinopathy Grading}
\author{\todo{authors and affiliations}}
\date{}

\begin{document}
\maketitle

\begin{abstract}
\textbf{Purpose.} Automated diabetic retinopathy (DR) grading must be accurate and must not
miss severe cases. \textbf{Methods.} TAG-MTL-DR couples EfficientNet-B3 lesion features
with a supervised vessel-segmentation branch, a topological vessel graph (end-points,
junctions and segments) encoded by a graph attention network, cross-attention from lesion
tokens to graph nodes, and a rank-consistent ordinal head. All methods were trained with
the same patient-grouped \NFolds-fold cross-validation over \NSeeds{} seeds and evaluated on
external data never used for training. \textbf{Results.} TAG-MTL-DR reached a QWK of
\OursQWK~$\pm$~\OursQWKsd{} (pooled \OursQWKpooled, 95\% CI \OursQWKlo--\OursQWKhi), against
\BestBaselineQWK{} for the best reproduced baseline (\BestBaseline; difference \DeltaBest,
Holm-adjusted $p$ = \PBest), with a severe non-referral rate of \OursSevNR\% (upper 95\%
bound \OursSevNRhi\%). On \ExtAName, QWK was \ExtAQWK{} and Sev-NR \ExtASevNR\%.
\textbf{Conclusions.} \todo{one sentence that the tables support, including the limits}
\end{abstract}

\section{Introduction}
\todo{Clinical need: screening volume, cost of missed severe cases. Gap: vessel topology
and the order of the grades are rarely modelled jointly with lesion features, and
comparisons rarely share one protocol. Contributions, stated as testable claims:
(i) the model; (ii) a single-protocol benchmark of reproduced baselines; (iii) ablations
isolating the vessel-path GRU, the lesion-aware nodes, the graph, the segmentation task, the
fusion and the ordinal head; (iv) external
validation with safety metrics and confidence intervals.}

\section{Related Work}
\todo{CNN and transformer graders; vessel segmentation; graph neural networks on vessel
structure; ordinal regression (CORAL). \Cref{tab:literature} lists results reported in the
literature; they are not comparable with ours because datasets, splits and metrics differ.}
\IfFileExists{\resultsdir/table_literature.tex}{\input{\resultsdir/table_literature.tex}}{}

\section{Methods}
\subsection{Model}
\textbf{Lesion branch (CNN).} EfficientNet-B3 feature map, one token per cell.
\textbf{Lesion-aware vessel graph (GNN).} The vessel skeleton is converted into a graph whose
nodes are end-points and junctions and whose edges are the vessel segments, weighted by
length and tortuosity. Each node carries its geometry (position, degree, type, distance to
the optic disc) \emph{and} the CNN feature sampled at its position, so that two graph
attention layers (GAT) propagate lesion evidence along the vessels.
\textbf{Vessel-path recurrence (RNN).} The vessel tree of each component is rooted at the
node closest to the optic disc; every root-to-leaf path is read in anatomical order
(disc to periphery) by a bidirectional GRU over the GAT node embeddings, and the paths are
pooled by attention.
\textbf{Fusion.} Multi-head cross-attention from the lesion tokens to the graph nodes and
path embeddings, with a residual connection; the result is concatenated with pooled lesion,
graph and path features. \textbf{Head.} CORAL with $K-1$ rank logits; temperature scaling
fitted on the validation patients. \textbf{Multi-task.} An auxiliary U-Net segments the
vessels (BCE + Dice) against annotated masks or classical pseudo-labels; loss
$\mathcal{L}_{\mathrm{CORAL}} + \lambda\,\mathcal{L}_{\mathrm{seg}}$, $\lambda = 0.5$.
\todo{architecture figure; state novelty only after a literature search.}

\subsection{Data and protocol}
\todo{datasets, grades, numbers of images and patients, external sets, ethics/licences.}
Splits are grouped by patient (both eyes in the same fold); 12.5\% of the training patients
form the validation set used for model selection (QWK). Every method uses the same folds,
seeds, input pipeline, class-balanced sampling, optimiser (AdamW, cosine schedule) and
early-stopping rule. External datasets are used only for evaluation, with the ensemble of
the cross-validated models.

\subsection{Metrics and statistics}
QWK, accuracy, macro-F1, AUC for referable DR (grade $\geq2$), severe and proliferative
non-referral rates (Sev-NR, PDR-NR: grade 3 or 4 predicted as 0 or 1) with exact
Clopper--Pearson intervals, Grade~0$\leftrightarrow$4 confusions, ECE and Brier score.
Confidence intervals: patient-level bootstrap. Comparison with each baseline: corrected
repeated $k$-fold $t$-test, Holm-adjusted; external sets: exact McNemar test and paired
bootstrap of the QWK difference. \todo{state the primary endpoint and the hypotheses
before running the experiments.}

\section{Results}
\input{\resultsdir/table_sota.tex}
\todo{comparison with reproduced baselines, \cref{tab:sota}.}
\IfFileExists{\resultsdir/table_ablation.tex}{\input{\resultsdir/table_ablation.tex}}{}
\todo{ablation, \cref{tab:ablation}: which component explains the gain.}
\IfFileExists{\resultsdir/table_external.tex}{\input{\resultsdir/table_external.tex}}{}
\todo{external validation, \cref{tab:external}.}
\input{\resultsdir/table_pergrade.tex}
\IfFileExists{\resultsdir/table_cost.tex}{\input{\resultsdir/table_cost.tex}}{}
\todo{computational cost, \cref{tab:cost}; data-efficiency curve (runs with
\texttt{--train\_fraction} 0.1, 0.25, 0.5, 1.0); qualitative explanations of successes
\emph{and} failures (\texttt{explain.py}).}

\begin{figure}[t]\centering
\includegraphics[width=0.48\linewidth]{\resultsdir/fig_confusion.pdf}\hfill
\includegraphics[width=0.40\linewidth]{\resultsdir/fig_reliability.pdf}
\caption{Normalised confusion matrix and reliability diagram of TAG-MTL-DR (pooled
out-of-fold predictions).}\label{fig:cm}
\end{figure}
\begin{figure}[t]\centering
\includegraphics[width=0.8\linewidth]{\resultsdir/fig_qwk.pdf}
\caption{Pooled QWK with 95\% patient-bootstrap intervals for all methods.}\label{fig:qwk}
\end{figure}

\section{Discussion}
\todo{what the results show; what they do not show; comparison with the literature only
where protocols match; failure cases (\cref{tab:pergrade}).}

\subsection{Limitations}
\todo{retrospective data; pseudo-label quality if no annotated masks; number of severe
cases and width of the Sev-NR interval; datasets without demographic data; no reader
study or prospective evaluation.}

\section{Conclusion}
\todo{one paragraph, consistent with the tables.}

\section*{Data and code availability}
Code: \todo{repository URL}. Datasets are public and obtained from their distributors.

\appendix
\section{Reporting checklist}
% Items expected by strong medical-imaging journals (inspired by CLAIM and TRIPOD+AI).
% Tick each item only when the manuscript actually contains it.
\begin{itemize}\small
\item[$\square$] Clinical question, intended use and target population stated.
\item[$\square$] Data sources, inclusion/exclusion criteria, image counts \emph{and} patient counts, grade distribution.
\item[$\square$] Reference standard: who graded, with which scale (ICDR), adjudication.
\item[$\square$] Splits at the patient level; no patient in two sets (asserted in \texttt{benchmark.py}).
\item[$\square$] External data never used for training, model selection or threshold tuning.
\item[$\square$] Primary endpoint and hypotheses fixed before the experiments.
\item[$\square$] All compared methods reproduced under one protocol; hyperparameters and budgets reported.
\item[$\square$] Several seeds; variability (SD) and confidence intervals reported.
\item[$\square$] Statistical tests suited to cross-validation (corrected $t$-test), multiplicity correction.
\item[$\square$] Safety metrics (Sev-NR, PDR-NR, 0$\leftrightarrow$4 confusions) with exact intervals.
\item[$\square$] Calibration (ECE, reliability diagram).
\item[$\square$] Ablation of every component claimed as a contribution.
\item[$\square$] Failure analysis (per-grade results, confusion matrix, examples).
\item[$\square$] Literature values shown separately, with sources and protocols, never mixed with reproduced results.
\item[$\square$] Limitations, including fairness and the absence of prospective evaluation.
\item[$\square$] Code, configuration and random seeds released.
\end{itemize}


\bibliographystyle{unsrt}
\bibliography{references}
\end{document}
